\documentclass[10pt,a4paper]{amsart}
\usepackage[margin=19mm]{geometry}
\usepackage{amsmath,amssymb,mathtools}
\usepackage[scr=rsfs,cal=euler]{mathalfa}
\usepackage{xfrac}
\usepackage[unicode,hidelinks,bookmarks=false]{hyperref}
\newcommand{\ie}{i.e.\ }
\newcommand{\cf}{cf.\ }
\newcommand{\resp}{respectively\ }
\newcommand{\Subsection}{\subsection}
\providecommand{\nobreakdash}{\nobreak\hbox{-}}
\setlength{\emergencystretch}{2em}
\multlinegap0pt
\usepackage{amssymb}
\usepackage{mathtools}
\usepackage{enumerate}
\usepackage{caption}
\newtheorem{theorem}{Theorem}
\newtheorem{proposition}{Proposition}
\newcommand\C{\mathbb{C}} 
\newcommand\Cscr{\mathscr{C}} 
\newcommand\D{\overline{\mathbb D}} 
\newcommand\Z{\mathbb{Z}} 
\newcommand\di{\partial} 
\newcommand\dibar{\overline\partial} 
\title{The nonhomogeneous Cauchy--Riemann equation on \\ families of open Riemann surfaces}
\author{Franc Forstneri{\v c}}
\date{Source: arXiv:2508.13660v1 (CC BY 4.0)}
\begin{document}
\maketitle

\noindent\emph{Source and licence.} The nonhomogeneous Cauchy--Riemann equation on \\ families of open Riemann surfaces. Version arXiv:2508.13660v1 (CC BY 4.0), distributed by arXiv under the Creative Commons Attribution 4.0 International licence, http://creativecommons.org/licenses/by/4.0/, as recorded in the arXiv licence metadata for this exact version and shown on the abstract page https://arxiv.org/abs/2508.13660v1. Full licence notice: \texttt{licenses/arxiv-2508.13660-CC-BY-4.0.txt}.\par\smallskip

\noindent\textit{Editorial scope.} Counted unit: the article's theorem th:dibar (source0.tex lines 517--536). The article's complete original proof of it is delivered with it (source0.tex lines 579--819, one \texttt{proof} environment of 241 lines, which the article places away from the statement). No mathematics is written, repaired or re-derived: the statement and the proof are the article's own lines, in the article's own notation, and no retained line is substituted or re-flowed.  Retained uncounted as the source-local context the proof needs: lines 498--500.  Nothing was omitted from any retained span, so the package carries no omission record.  The retained lines cite 3 works, whose entries are transcribed verbatim from the article's own bibliography source in the e-print and delivered with short editorial labels recorded in the item's reference mapping.  No retained line contains a figure environment, an image-inclusion command or drawing code, so no image is delivered and the package declares no asset dependency.  Editorial formatting only, in addition: an AMS \texttt{amsart} preamble that restates the article's own declarations and macros used by the retained lines, namely the environments \texttt{theorem}, \texttt{proposition} on separate counters, together with the standard packages those lines need.  The retained environments print in this document as Theorem 1, Proposition 1 in order of appearance, whereas the article prints the same items with the numbering of its own full document.  The complete native source of the article as distributed in the arXiv e-print for this exact version is bundled unchanged as \texttt{sources/183496/source0.tex}.
\begin{equation}\label{eq:derivatives}
	df=\di_J f+\dibar_J f=f_z \,\cdotp dz + f_{\bar z} \,\cdotp d\bar z.  
\end{equation}

\begin{theorem}\label{th:dibar}
Assume that $X$ is a smooth open orientable surface, 
$\Omega\Subset X$ is a relatively compact domain with 
$\Cscr^{(k+1,\alpha)}$ boundary for some $k\in\Z_+$ and
$0<\alpha<1$, $l\in \Z_+$, $B$ is a paracompact Hausdorff
space if $l=0$ and a $\Cscr^l$ manifold if $l>0$,
and $\{J_b\}_{b\in B}$ is a family of complex structures 
of class $\Cscr^{l,(k,\alpha)}(B\times \overline\Omega)$ 
on $\overline\Omega$.
Given a family of $(0,1)$-forms $\{\beta_b\}_{b\in B}$ 
% \subset \Gamma(\overline\Omega,T^{*(0,1)}_{J_b}\overline\Omega)$ 
on $\overline\Omega$ of class 
$\Cscr^{l,(k,\alpha)}(B\times \overline\Omega)$, 
there is a function 
$f\in \Cscr^{l,(k+1,\alpha)}(B\times\overline\Omega)$ satisfying
\begin{equation}\label{eq:dibar}
	\dibar_{J_b} f(b,\cdotp) = \beta_b \ \ 
	\text{on $\overline \Omega$ for every $b\in B$.}
\end{equation}
\end{theorem}

\begin{proof}[Proof of Theorem \ref{th:dibar}]
We shall use the connection between complex
structures and Beltrami multipliers; see 
\cite[Sect.\ 2]{Forstneric2024Runge} or any 
standard text on quasiconformal maps for the details.

Choose a smooth complex structure $J$ on $X$ and 
a $J$-holomorphic immersion $z:X\to\C$.
Let $\Omega\Subset X$ be as in the theorem. 
Every function $\mu:\overline \Omega\to\D$ 
of class $\Cscr^{(k,\alpha)}(\overline \Omega)$  with values
in the unit disc $\D$ 
determines a complex structure $J_\mu$ on $\overline\Omega$ 
of the same class, with $\mu=0$ corresponding to $J|_{\overline\Omega}$. 
Conversely, every complex structure $J'$ on $\overline \Omega$
of class $\Cscr^{(k,\alpha)}$ in the same orientation class 
as $J$ equals $J_\mu$ for a unique 
$\mu\in \Cscr^{(k,\alpha)}(\overline \Omega,\D)$. 
A differentiable function $f:\overline\Omega\to\C$ satisfies 
$\dibar_{J_\mu}f=0$, and hence is $J_\mu$-holomorphic on
$\Omega$, if and only if it satisfies the homogeneous Beltrami equation 
$f_{\bar z} = \mu f_z$, where the partial derivatives $f_z$ and 
$f_{\bar z}$ are defined by \eqref{eq:derivatives}. The analogous
statements hold for globally defined complex structures on $X$
and function $\mu:X\to \D$.

To prove the theorem, it suffices to show that for every point $b_0\in B$ 
there is a neighbourhood $B_0\subset B$ of $b_0$ such that 
equation \eqref{eq:dibar} is solvable on $B_0\times \overline\Omega$ with 
$f\in \Cscr^{l,(k+1,\alpha)}(B_0\times \overline\Omega)$. 
By using $\Cscr^l$ partitions of unity on $B$ we then obtain 
a solution on $B\times \overline \Omega$, thereby proving the theorem.

Note that the complex structure $J_{b_0}$ on $\overline\Omega$
extends to a complex structure on $X$ of local class $\Cscr^{(k,\alpha)}$.
To see this, we represent 
$J_{b_0}$ in terms of $J$ by a Beltrami multiplier 
$\mu\in \Cscr^{(k,\alpha)}(\overline \Omega,\D)$. 
Since $b\Omega$ is of class $\Cscr^{(k+1,\alpha)}$,
$\mu$ extends from $\overline \Omega$ to a function
$\mu':X\to \D$ of class $\Cscr^{(k,\alpha)}$ 
with compact support 
(see Gilbarg and Trudinger \cite[Lemma 6.37]{GilbargTrudinger1983}). 
The associated complex structure $J_{\mu'}$ on $X$ is of local 
class $\Cscr^{(k,\alpha)}$ and it coincides with $J_\mu$ on 
$\overline \Omega$. For $k\ge 1$ the same conclusion 
holds if $b\Omega$ is of class $\Cscr^{(k,\alpha)}$.

This reduces the proof to the following proposition.
In this result, the smooth structure on $X$ is induced  
by the complex structure $J$, and the H\"older norms are 
computed with respect to a fixed smooth Riemannian metric on $X$
whose precise choice is unimportant. 

%
%  Solving the dibar equation
%
\begin{proposition}\label{prop:dibar}
Let $(X,J)$ be an open Riemann surface, 
and let $\Omega\Subset X$, $k$, and $\alpha$ be as in Theorem \ref{th:dibar}.  
For $\mu\in \Cscr^{(k,\alpha)}(\overline \Omega,\D)$ let $J_\mu$ 
denote the associated complex structure on $\overline\Omega$, 
with $J_0=J|_{\overline \Omega}$. For $c>0$ set 
$
	B_c=\{\mu \in \Cscr^{(k,\alpha)}(\overline\Omega) : 
	\|\mu\|_{k,\alpha} <c\}.
$
There exists $c>0$ such that for any map 
$B_c\ni \mu \mapsto \beta_\mu \in  
\Gamma^{(k,\alpha)}(\overline\Omega,T^{*(0,1)}_{J_\mu}\overline\Omega)$
of class $\Cscr^l$, $l\in \{0,1,\ldots,\infty,\omega\}$, 
there is a function 
$f\in \Cscr^{l,(k+1,\alpha)}(B_c\times \overline\Omega)$ 
such that for every $\mu \in B_c$
the function $f_\mu=f(\mu,\cdotp):\overline\Omega\to\C$ satisfies 
\begin{equation}\label{eq:dibar2}
	\dibar_{J_\mu} f_\mu = \beta_\mu. 
\end{equation}
\end{proposition}

\begin{proof}
We begin by recalling some
technical tools from \cite[Secs.\ 3-4]{Forstneric2024Runge}.
Choose a $J$-holomorphic immersion $z:X\to\C$.  
There is a Cauchy kernel on $(X,J)$ which determines on any 
smoothly bounded domain $\Omega\Subset X$ a pair 
of bounded linear operators 
$P:\Cscr^{(k,\alpha)}(\overline \Omega) \to \Cscr^{(k+1,\alpha)}(\overline \Omega)$ and 
$S:\Cscr^{(k,\alpha)}(\overline \Omega) \to \Cscr^{(k,\alpha)}(\overline \Omega)$
($k\in\Z_+$, $0<\alpha<1$) such that 
for any $\phi\in \Cscr^{(k,\alpha)}(\overline \Omega)$, the 
Cauchy operator $P$ solves the equation
$P(\phi)_{\bar z}=\phi$ on $\overline \Omega$ and 
the Beurling operator $S$ is given by $S(\phi)=P(\phi)_z$. 
The properties of these operators are summarised in 
\cite[Theorem 3.2]{Forstneric2024Runge}. 
Although the cited result is stated for domains $\Omega$
with $\Cscr^\infty$ boundaries, it is clear from the proof and 
\cite[Lemma 6.37]{GilbargTrudinger1983} that it holds 
if $b\Omega$ is of class $\Cscr^{(k,\alpha)}$ if $k\ge 1$, 
and of class $\Cscr^{(1,\alpha)}$ if $k=0$.

By \cite[Theorem 4.1]{Forstneric2024Runge} there are a constant 
$c>0$ and a function $h:B_c \times \overline \Omega\to \C$ such that
for every $\mu\in B_c$, $h_\mu=h(\mu,\cdotp):\overline \Omega\to\C$ 
is a $J_\mu$-holomorphic immersion of class 
$\Cscr^{(k+1,\alpha)}(\overline \Omega)$ depending analytically on $\mu$.
We recall the proof since we shall use a similar idea 
in the sequel. The function $f=h_\mu$ must solve 
the Beltrami equation $f_{\bar z} = \mu f_z$.
We look for a solution in the form $f=z|_{\overline \Omega} +P(\phi)$ with
$\phi\in \Cscr^{(k,\alpha)}(\overline \Omega)$.
Note that $\phi=0$ corresponds to $f=z|_{\overline\Omega}$. 
We have that 
\[
	f_{\bar z}= P(\phi)_{\bar z}  = \phi
	\quad \text{and} \quad 
	f_z = 1 + P(\phi)_z = 1+ S(\phi).
\]
Inserting in the Beltrami equation $f_{\bar z} = \mu f_z$ gives $(I - \mu S) \phi =\mu$,
where $I$ denotes the identity map on $\Cscr^{(k,\alpha)}(\overline\Omega)$. 
For $\|\mu\|_{k,\alpha}$ small enough we have that $\|\mu S\|_{k,\alpha}<1$, 
so the operator $I- \mu S$ is invertible and its bounded inverse  
depends analytically on $\mu$: 
\[ 
	\Theta(\mu) := (I - \mu S)^{-1} =
	\sum_{j=0}^\infty (\mu S)^j 
	\in  \mathrm{Lin}(\Cscr^{k,\alpha}(\overline\Omega)). 
\] 
(See Mujica \cite[29.3 Theorem]{Mujica1986}.)
Hence, $\phi=\Theta(\mu)\mu$ and we get 
the following solution $h_\mu=f$ to the 
Beltrami equation $f_{\bar z} = \mu f_z$ on $\overline\Omega$:
\[ 
	h_\mu = z|_{\overline\Omega} +  P (\Theta(\mu) \mu) 
	= z|_{\overline\Omega} + P((I - \mu S)^{-1} \mu)
	 \in \Cscr^{(k+1,\alpha)}(\overline\Omega).
\] 
Note that $h_\mu$ depends analytically on $\mu$. For $\|\mu\|_{k,\alpha}$ 
small enough, $h_\mu$ is so close to $h_0=z|_{\overline\Omega}$ in
$\Cscr^{(k+1,\alpha)}(\overline\Omega)$ that it is an immersion.
Hence, for $c>0$ small enough, 
$\{\theta_\mu=dh_\mu\}_{\mu\in B_c}$
is a family of nowhere vanishing holomorphic 
$1$-forms on $(\Omega,J_b)$ of class 
$\Cscr^{(k,\alpha)}(\overline\Omega)$ with analytic dependence on $\mu$. 
(The existence of a family $\{\theta_b\}_{b\in B}$ of globally defined nowhere vanishing holomorphic $1$-forms on $(X,J_b)$ of class $\Cscr^{l,(k,\alpha)}$ was proved in \cite[Theorem 7.1]{Forstneric2024Runge} under a stronger assumption on the parameter space $B$.) 
%
The conjugate $\bar \theta_\mu=d\overline{h_\mu}$ is a nowhere vanishing 
antiholomorphic $(0,1)$-form with respect to $J_\mu$ for every 
$\mu\in B_c$. Thus, every family $\{\beta_\mu\}_{\mu\in B_c}$ 
on $\overline\Omega$, where $\beta_\mu$ is a $(0,1)$-form 
with respect to $J_\mu$, is of the form 
$\beta_\mu=u_{\mu}\, d\overline{h_\mu}$ for a family of functions 
$u_\mu: \overline\Omega\to \C$, $\mu\in B_c$. If the family 
$\{\beta_\mu\}_{\mu\in B_c}$ is of class 
$\Cscr^{l,(k,\alpha)}(B_c\times\overline\Omega)$
then $\{u_\mu\}_{\mu\in B_c}$ is of the same class, and vice versa.

We shall express the $\dibar_{J_\mu}$-equation \eqref{eq:dibar2} as a 
nonhomogeneous Beltrami equation with respect to the immersion $z$.
For $\mu\in B_c$ we can uniquely express any complex $1$-form 
$\beta$ on $\overline\Omega$ as 
\begin{equation}\label{eq:beta}
	\beta=A dz+Bd\bar z = A_\mu dh_\mu + B_\mu d\overline{h_\mu}.
\end{equation}
Note that $\beta$ is of class $\Cscr^{(k,\alpha)}(\overline\Omega)$
if and only if the coefficients $A,B,A_\mu,B_\mu:\overline\Omega\to\C$
are of this class. We now express $A_\mu$ and $B_\mu$ in terms of 
the functions $A,B,\mu$, and 
\begin{equation}\label{eq:gmu}	
	g_\mu := (h_\mu)_z \in \Cscr^{(k,\alpha)}(\overline\Omega).
\end{equation}
We have that 
\[
	dh_\mu = (h_\mu)_z \, dz + (h_\mu)_{\bar z} \, d\bar z
	= g_\mu\, dz + \mu g_\mu \, d\bar z
\]
where the second identity follows from the Beltrami equation
$(h_\mu)_{\bar z}=\mu (h_\mu)_z$. From this and \eqref{eq:beta} it follows that
\begin{eqnarray*}
	A dz+Bd\bar z &=& A_\mu (g_\mu dz + \mu g_\mu d\bar z)
	+ B_\mu (\overline {\mu g_\mu}\, dz + \overline{g_\mu} \, d\bar z) \\
	&=& (A_\mu g_\mu + B_\mu \overline{\mu g_\mu}) dz
	+ (A_\mu \mu g_\mu + B_\mu \overline{g_\mu}) d\bar z
\end{eqnarray*}
and hence
\[
	A = A_\mu g_\mu + B_\mu \overline{\mu g_\mu},
	\qquad
	B = A_\mu \mu g_\mu + B_\mu \overline{g_\mu}.
\]
Solving these equations on $A_\mu$ and $B_\mu$ gives
\[
	A_\mu = \frac{A  - \bar \mu B}{(1-|\mu|^2)g_\mu}, %(A  - \bar \mu B),
	\qquad 
	B_\mu = \frac{B  - \mu A}{(1-|\mu|^2) \overline{g_\mu}}.
\]
Taking $\beta=df$ we have $A=f_z$, $B=f_{\bar z}$, $A_\mu = f_{h_\mu}$, 
and $B_\mu= f_{\overline{h_\mu}}$. Inserting these quantities in 
the above expression for $B_\mu$ shows that the 
nonhomogeneous Cauchy--Riemann equation 
\[
	\dibar_{J_\mu} f = \beta_\mu = u_\mu d\overline{h_\mu} 
	\ \Longleftrightarrow \ 
	f_{\overline{h_\mu}} = 
	\frac{f_{\bar z}-\mu f_z} {(1-|\mu|^2) \overline{g_\mu}} = u_\mu
\]
is equivalent to the nonhomogeneous Beltrami equation 
\begin{equation}\label{eq:dibarmu}
	f_{\bar z}-\mu f_z = (1-|\mu|^2) \overline{g_\mu} \, u_\mu
\end{equation}
with $g_\mu$ given by \eqref{eq:gmu}. 
The regularity conditions in the proposition imply that the right hand side is of class $\Cscr^{l,(k,\alpha)}$ on $B_c\times \overline\Omega$. 
We look for a solution of \eqref{eq:dibarmu} in the form $f=f(\mu)=P(\phi)$ with 
$\phi\in \Cscr^{(k,\alpha)}(\overline\Omega)$ to be determined. 
Inserting $f_{\bar z}=P(\phi)_{\bar z}=\phi$ and $f_{z}=P(\phi)_z=S(\phi)$
into \eqref{eq:dibarmu} gives
\[
	f_{\bar z}-\mu f_z = (I-\mu S)\phi = 
	(1-|\mu|^2) \overline{g_\mu} \, u_\mu.
\]
For $\|\mu\|_{k,\alpha}$ small enough the operator $I-\mu S$ is
invertible and we obtain 
\[
	\phi=\phi_\mu = 
	(I-\mu S)^{-1}\left((1-|\mu|^2) \overline{g_\mu} \, u_\mu \right).
\] 
Since the bounded linear operator 
$(I-\mu S)^{-1}\in \mathrm{Lin}(\Cscr^{k,\alpha}(\overline \Omega))$
is analytic in $\mu$ and $(1-|\mu|^2) \overline{g_\mu} \, u_\mu 
\in \Cscr^{l,(k,\alpha)}(B_c\times\overline \Omega)$, the map
$(\mu,x)\mapsto \phi_\mu(x)$ belongs to
$\Cscr^{l,(k,\alpha)}(B_c\times\overline \Omega)$.
Finally, the solution of \eqref{eq:dibarmu} is $f_\mu = P(\phi_\mu)$,  
and the map $(\mu,x)\to f_\mu(x)$ belongs to
$\Cscr^{l,(k+1,\alpha)}(B_c\times \overline \Omega)$. 
\end{proof}

This completes the proof of Theorem \ref{th:dibar}. 
\end{proof}

\begin{thebibliography}{99}
\bibitem[Forstn]{Forstneric2024Runge} F.~Forstneri\v{c}. \newblock Runge and {M}ergelyan theorems on families of open {R}iemann surfaces. \newblock {\em arXiv e-prints}, 2024. \newblock \url{https://arxiv.org/abs/2412.04608}.
\bibitem[Gilbar]{GilbargTrudinger1983} D.~Gilbarg and N.~S. Trudinger. \newblock {\em Elliptic partial differential equations of second order. 2nd ed}, volume 224 of {\em Grundlehren Math. Wiss.} \newblock Springer, Cham, 1983.
\bibitem[Mujica]{Mujica1986} J.~Mujica. \newblock {\em Complex analysis in {Banach} spaces. {Holomorphic} functions and domains of holomorphy in finite and infinite dimensions}, volume 120 of {\em North-Holland Math. Stud.} \newblock Elsevier, Amsterdam, 1986.
\end{thebibliography}
\end{document}
